\documentclass[10pt,a4paper]{amsart}
\usepackage[margin=19mm]{geometry}
\usepackage{amsmath,amssymb,mathtools}
\usepackage[scr=rsfs,cal=euler]{mathalfa}
\usepackage{xfrac}
\usepackage[unicode,hidelinks,bookmarks=false]{hyperref}
\newcommand{\ie}{i.e.\ }
\newcommand{\cf}{cf.\ }
\newcommand{\resp}{respectively\ }
\newcommand{\Subsection}{\subsection}
\providecommand{\nobreakdash}{\nobreak\hbox{-}}
\setlength{\emergencystretch}{2em}
\multlinegap0pt
\usepackage{amssymb}
\newtheorem{lemma}{Lemma}
\newcommand{\eps}{\varepsilon}
\renewcommand{\ge}{\geqslant}
\renewcommand{\le}{\leqslant}
\title{Maximal Kolmogorov Complexity in a Hamming Ball}
\author{Alexander Kozachinskiy, Nikolay Vereshchagin}
\date{Source: arXiv:2609.11362v1 (CC BY 4.0)}
\begin{document}
\maketitle

\noindent\emph{Source and licence.} Maximal Kolmogorov Complexity in a Hamming Ball. Version arXiv:2609.11362v1 (CC BY 4.0), distributed by arXiv under the Creative Commons Attribution 4.0 International licence, http://creativecommons.org/licenses/by/4.0/, as recorded in the arXiv licence metadata for this exact version and shown on the abstract page https://arxiv.org/abs/2609.11362v1. Full licence notice: \texttt{licenses/arxiv-2609.11362-CC-BY-4.0.txt}.\par\smallskip

\noindent\textit{Editorial scope.} Counted unit: the article's lemma lem:construct (source0.tex lines 457--462). The article's complete original proof of it is delivered with it (source0.tex lines 464--552, one \texttt{proof} environment of 89 lines). No mathematics is written, repaired or re-derived: the statement and the proof are the article's own lines, in the article's own notation, and no retained line is substituted or re-flowed.  Retained uncounted as the source-local context the proof needs: lines 227--229; lines 232--236; lines 421--425.  Nothing was omitted from any retained span, so the package carries no omission record.  No retained line contains a citation, so the wrapper carries no bibliography and no bibliography entry.  No retained line contains a figure environment, an image-inclusion command or drawing code, so no image is delivered and the package declares no asset dependency.  Editorial formatting only, in addition: an AMS \texttt{amsart} preamble that restates the article's own declarations and macros used by the retained lines, namely the environments \texttt{lemma} on separate counters, together with the standard packages those lines need.  The retained environments print in this document as Lemma 1 in order of appearance, whereas the article prints the same items with the numbering of its own full document.  The complete native source of the article as distributed in the arXiv e-print for this exact version is bundled unchanged as \texttt{sources/210011/source0.tex}.
\begin{equation}\label{eq:harper}
  |A|\ \ge\ V(m)\quad\Longrightarrow\quad |B_a(A)|\ \ge\ V(m+a).
\end{equation}

\begin{equation}\label{eq:ratio}
  \frac{V(s)}{V(s-1)}\le n+1\, \text{ for all }s\ge1,
  \qquad
  \frac{V(s)}{V(s-1)}\ge1+\frac1n\ \text{ for }1\le s\le n/2 .
\end{equation}

\begin{lemma}[centers with small covering multiplicity]\label{lem:pack}
Let $s\ge0$ be any radius and let $M$ be such that
$M\cdot V(s)\le 2n\cdot2^n$. Then there exist $x_1,\dots,x_M\in\{0,1\}^n$ such that every $y\in\{0,1\}^n$ belongs to at
most $8n$ of the balls $B_s(x_i)$.
\end{lemma}

\begin{lemma}\label{lem:construct}
Let $q\le n/2$ be such that $V(q)\le2^k$. Set  $M:=\lfloor2^k/V(q)\rfloor\ge1$.
Assume further that 
$M\cdot V(q+r)\le2n\cdot2^n$. Then there is a
$\bigl(k,\,r,\,\lceil\log (M V(q+r))\rceil,\,1/32n^2\bigr)$-concentrator.
\end{lemma}

\begin{proof}
Let $x_1,\dots,x_M$ be as in Lemma~\ref{lem:pack} with $s=q+r$: every string is
covered by at most $8n$ of the balls $B_{q+r}(x_i)$; in particular it is covered
by at most $8n$ of the smaller balls $B_{q}(x_i)$. Put
\[
  C\ =\ \bigcup_{i=1}^{M}B_q(x_i),\qquad\text{so}\qquad
  B_r(C)\ =\ \bigcup_{i=1}^{M}B_{q+r}(x_i).
\]

\emph{Sizes.} Counting with multiplicities, $$MV(q)\ge|C|\ge MV(q)/(8n)$$ and
$$MV(q+r)\ge|B_r(C)|\ge MV(q+r)/(8n).$$ 
Here $MV(q)$ is close to $2^k$  by the choice of $M$. Indeed, $MV(q)\le2^k$ and
$$MV(q)\ge2^k-V(q)\ge2^{k-1}$$ (if $V(q)\le2^{k-1}$ this is clear, and otherwise
$M=1$ and $MV(q)=V(q)>2^{k-1}$). With $l:=\lceil\log(MV(q+r))\rceil$,
conditions (C1) and (C2) hold with $\eps\le1/16n$. Indeed
$$
|C|\le MV(q)\le 2^k \text{ and } 
|C|\ge MV(q)/8n\ge2^{k-1}/8n.
$$ 
Besides,  $2^{l-1}<MV(q+r)\le 2^l$ by the choice of $l$. Hence 
$$
|B_r(C)|\le MV(q+r)\le 2^l \text{ and } |B_r(C)|\ge MV(q+r)/8n>2^{l-1}/8n.
$$

\emph{Expansion.} Let $\eps=1/32n^2$ and let $C'\subseteq C$ satisfy
$|C'|\ge(1-\eps)|C|$; we have to bound $|B_r(C')|$ from below.

\emph{Step 1: most of the balls $B_q(x_i)$ lose few elements.} Consider the quantity
\[
  \Sigma\ :=\ \sum_{i=1}^{M}\bigl|B_q(x_i)\setminus C'\bigr| ,
\]
accounting for the total loss of all the balls, where 
each lost string $y$ is counted once for every index $i$ with $y\in B_q(x_i)$
--- so a string lying in several of the balls is counted several times.
 Every string counted here lies in
$C\setminus C'$, because $B_q(x_i)\subseteq C$, and it lies in at most $8n$ of
the balls $B_q(x_i)$. Hence
\[
  \Sigma\ =\ \sum_{y\in C\setminus C'}\#\{i:\ y\in B_q(x_i)\}\ \le\
  8n\,|C\setminus C'|\ \le\ 8n\eps|C|\ \le\ 8n\eps MV(q),
\]
the last step because $|C|\le MV(q)$. Now $\Sigma$ is a sum of $M$ non-negative
terms, so fewer than $M/2$ of them exceed $2\Sigma/M$; consequently for at
least $M/2$ indices $i$ (call them \emph{good})
\[
  |B_q(x_i)\setminus C'|\ \le\ \frac{2\Sigma}{M}\ \le\ 16n\eps V(q)\ =\
  \frac{V(q)}{2n},
\]
the last step because  $\eps=1/32n^2$.

\emph{Step 2: a good ball still contains more than $V(q-1)$ elements.} Since
$|B_q(x_i)|=V(q)$, for a good $i$ we get
\[
  |B_q(x_i)\cap C'|\ =\ V(q)-|B_q(x_i)\setminus C'|\ \ge\
  V(q)\Bigl(1-\frac1{2n}\Bigr)\ >\ V(q-1) .
\]
The last inequality is trivial for $q=0$ (then $V(q-1)=0$), and for $q\ge1$ it
holds because $V(q)\ge(1+1/n)V(q-1)$ by the second inequality
of~\eqref{eq:ratio} --- this is the only place where the hypothesis $q\le n/2$
is used --- so that $V(q-1)\le\frac n{n+1}V(q)$, while
$1-\frac1{2n}>\frac n{n+1}$ for $n\ge2$.

\emph{Step 3: Harper.} Applying~\eqref{eq:harper} to the set
$B_q(x_i)\cap C'$, whose cardinality is at least that of a ball of radius
$q-1$, we get for every good $i$
\[
  \bigl|B_r\bigl(B_q(x_i)\cap C'\bigr)\bigr|\ \ge\ V(q-1+r)\ \ge\
  \frac{V(q+r)}{n+1},
\]
the last step by the first inequality of~\eqref{eq:ratio}, which is valid for
all radii --- note that $q+r$ may well exceed $n/2$ here.

\emph{Step 4: from multiplicities back to the union.} For a good $i$ the set
$B_r(B_q(x_i)\cap C')$ is contained in $B_r(C')$, and also in $B_{q+r}(x_i)$.
Every string belongs to at most $8n$ of the balls $B_{q+r}(x_i)$ and hence to
at most $8n$ of these sets, so the union of the latter, which is a subset of
$B_r(C')$, is at least their total size divided by $8n$:
\[
  |B_r(C')|\ \ge\ \frac1{8n}\sum_{i\ \text{good}}
  \bigl|B_r\bigl(B_q(x_i)\cap C'\bigr)\bigr|\ \ge\
  \frac1{8n}\cdot\frac M2\cdot\frac{V(q+r)}{n+1}\ \ge\
  \frac{MV(q+r)}{32n^2} .
\]
Finally $|B_r(C)|\le MV(q+r)$ and $\eps=1/32n^2$, so
$|B_r(C')|\ge\eps\,|B_r(C)|$. This is (C3). 
%(Steps~1 and~4 are what fix the
%value of $\eps$: both hold with equality for $\eps=1/32n^2$, while (C1) and
%(C2) need only $\eps\le1/16n$.)
\end{proof}

\end{document}
